% 復習2　問いの立て方
\session{2}{1}{60分}{問いの立て方}{AIとブレスト}{}

\goals{
  \item テーマとリサーチ・クエスチョン（問い）の違いを説明できる
  \item 良い問いの条件（調べて答えられる／範囲が適切／答える意義がある）を挙げられる
  \item AIとのブレストで視点を広げつつ、最後は自分で選ぶ
}

\work{確認}{テーマと問いの違い}
\begin{wstab}{colspec={Q[l,wd=40mm,m] X[l,m]},row{1-Z}={ht=9mm}}
  \hc{例：テーマ} & 小売業のキャッシュレス決済 \\
  \hc{例：問い} & なぜ小規模な飲食店では、キャッシュレス決済の導入が遅れているのか \\
  \hc{自分のテーマ} & \\
  \hc{問いにすると} & \\
  \hc{良い問いの3条件} & ①\hspace{38mm} ②\hspace{38mm} ③ \\
\end{wstab}

\work[60mm]{ワーク1}{AIで視点を広げる}
\begin{promptbox}[視点を広げる]
私は大学生で、「（あなたのテーマ）」に関心があります。このテーマについて、経営戦略、組織、マーケティング、情報システム、経済、社会の6つの視点から、それぞれ2つずつ研究の問いを提案してください。各問いは「なぜ〜なのか」「〜は〜にどのように影響するのか」の形で、1行ずつ箇条書きにしてください。
\end{promptbox}
\instr{AIの候補のうち、気になったものを視点ごとに書き写す（言い換えてよい）。}
\begin{wstabh}{colspec={Q[l,wd=22mm,m] X[l,m]},row{2-Z}={ht=10mm}}
  視点 & 気になった問いの候補 \\
  経営戦略 & \\ 組織 & \\ マーケティング & \\ 情報システム & \\ 経済 & \\ 社会 & \\
\end{wstabh}

\work[60mm]{ワーク2}{3条件で採点し、上位3つに絞る}
\instr{各条件を 3（十分）・2（やや）・1（不十分）で採点する。必要ならAIに評価と狭める案を求める。}
\begin{promptbox}[問いを絞る]
次の問いについて、「大学生が数か月で調べて答えられるか」「範囲は適切か」「答えることに意義があるか」の3点から評価し、必要なら狭める案を示してください。問い：（候補の問い）
\end{promptbox}
\begin{wstabh}{colspec={X[l,m] Q[c,wd=17mm,m] Q[c,wd=17mm,m] Q[c,wd=17mm,m] Q[c,wd=12mm,m] Q[c,wd=10mm,m]},row{2-Z}={ht=10mm}}
  問いの候補 & 調べて\newline 答えられる & 範囲が\newline 適切 & 答える\newline 意義 & 合計 & 順位 \\
  & & & & & \\ & & & & & \\ & & & & & \\ & & & & & \\ & & & & & \\
\end{wstabh}

\work{セルフチェック}{「それはどうやって調べるの？」に答える}
\instr{上位の問いについて、AIに「この問いを調べる方法について、学生に鋭い質問を3つしてください」と頼み、答えてみる。友人と説明し合ってもよい。}
\twoboxes{受けた質問}{自分の答え・気づいたこと}{30mm}

\work{決定}{仮のリサーチ・クエスチョン}
\answerbox{18mm}

\begin{nexttime}
  \item 仮の問いを選んだ理由を200字程度でまとめる
  \item 使ったプロンプトと、AIの案のうち採用しなかったものとその理由を、下のAI活用記録に書く
\end{nexttime}
\answerbox[選んだ理由（200字程度）]{40mm}
\aimemo
